\section{第\ref{sec:synthesis}章　マシン全体}

まとめの章なので新しい実験は挙げない。各行は、その主張を詳しく扱った章と同じ出典にもとづく。\nt{GLUT}バスと\nt{GABA}による流れの制御は確立しているが、ブロードキャストチャネルの役割とその協調は大部分が仮説である。

{\footnotesize
\setlength{\tabcolsep}{3pt}\renewcommand{\arraystretch}{1.15}
\begin{longtable}{>{\raggedright}p{0.5cm}>{\raggedright}p{4.0cm}>{\raggedright}p{5.6cm}>{\raggedright}p{2.3cm}>{\raggedright\arraybackslash}p{1.7cm}}
\toprule
No. & 主張 & 実験と測定されたもの & 出典 & 状態 \\
\midrule\endfirsthead
\toprule
No. & 主張 & 実験と測定されたもの & 出典 & 状態 \\
\midrule\endhead
1 & $8.6\cdot10^{10}$ 個のニューロンに対し、制御信号はわずか四種類~\eqref{eq:machine} & 脳のホモジネート中の細胞核の計数（等方性分画法）：成人男性でニューロンは $(86.1\pm8.1)\cdot10^9$ 個 & \cite{Azevedo2009} & 測定済み \\
2 & 命題~\ref{prop:architecture}の最初の二層：データは\nt{GLUT}のアドレスバスを流れ、結合の符号は送り手が決め、流れは\nt{GABA}ニューロンが方向づけ正規化する（第\ref{sec:neurontypes}, \ref{sec:gaba}章） & ニューロンあたり主要な伝達物質は一つ（測定の総説）；フィードフォワード抑制は錐体ニューロンが入力を加算する窓を狭める；全体の活動による除算は多くの領域で測定されている & \cite{Strata1999,Pouille2001,Carandini2012} & 測定済み \\
3 & 素子は信頼できない：シナプスはスパイクの大部分を失う（表~\ref{tab:computer}、第\ref{sec:code}章） & 海馬スライス、最小刺激：スパイクの半分以上がCA1で応答を起こさない。閾値や軸索伝導の失敗は除外され、原因は確率的な伝達物質放出 & \cite{AllenStevens1994} & 測定済み \\
4 & 脳は $\sim20$ ワットで動き、ニューロンあたり平均で毎秒約1スパイク（第\ref{sec:code}章）；このようなマシンをエンジニアはまだ作れていない & 測定されたコストからの推定~\eqref{eq:energy}：シナプスの仕事を含め1スパイクあたり約 $10^9$ 個のATP分子（げっ歯類の皮質）；皮質についての詳細な推定ではさらに低い発火率になる。技術の現状は総説による & \cite{Attwell2001,Lennie2003,MehonicKenyon2022} & 理論 \\
5 & 大多数のシナプスの学習は局所的な適格度トレースと一つの共通の数 $\delta$ の積（\eqref{eq:machine}の第2式、第\ref{sec:dofa}章） & サルの\nt{DOPA}ニューロンは報酬予測誤差に応答する；線条体スライスでは、ドーパミンがグルタミン酸入力の $0.3$--$2$\,s 後に来たときだけスパインを大きくする。トレースが測定された & \cite{Schultz1997,Yagishita2014} & 測定済み \\
6 & 出力ごとの専用の誤差配線があるのは運動学習の回路だけ（第\ref{sec:motorlearning}章） & 小脳：登上線維の信号と入力が同時に来るとその入力が弱まる；サルではプルキンエ細胞の複雑スパイクが運動誤差の方向を運び、補正を起こす & \cite{Ito2001,Herzfeld2018} & 測定済み \\
7 & 各ブロードキャストチャネルは数本の線の束（命題~\ref{prop:bundle}） & \nt{DOPA}について：同じ課題で異なるニューロンが報酬、運動、刺激の特徴を運ぶ；速い誤差と遅いドーパミンレベルは別物。\nt{SERO}, \nt{NORA}, \nt{ACET}ではこの分解はあまり調べられていない & \cite{Engelhard2019,Hamid2016,Mohebi2019} & 測定済み \\
8 & \nt{SERO}：レベル調整器であり、仮説では時間幅 $\tau_\gamma$ も決める（第\ref{sec:serocomputer}章） & 自己抑制：自らの5-HT$_{1A}$受容体のアゴニストが縫線核のセロトニンニューロンを抑制する（測定済み）。時間幅は理論で提案された；最も強い実験は、マウスでこれらのニューロンを光遺伝学的に活性化すると遅延報酬を待つ時間が延びること & \cite{Sprouse1987,Doya2002,Miyazaki2014} & 仮説 \\
9 & \nt{NORA}：ゲイン $g$ が探索か決定かを選ぶ（第\ref{sec:nora}章） & SN比の向上：覚醒したサルの聴覚皮質で、ノルアドレナリンは背景活動を同種の鳴き声への応答より強く抑える（測定済み）；乗算的ゲインはモデル；探索と決定の選択での役割は理論 & \cite{Foote1975NE,ServanSchreiber1990,AstonJones2005} & 仮説 \\
10 & \nt{ACET}：書き込みと読み出しを切り替える（第\ref{sec:acet}章） & 三つの作用（内部結合の弱化、入力の強化、可塑性の促進）は測定済み；モードの切り替えはヒトでのスコポラミン実験が間接的に支持 & \cite{HasselmoBower1992,Gil1997,Atri2004} & 仮説 \\
11 & 式~\eqref{eq:machine}はマシン全体を記述する & 第\ref{sec:glutcomputer}--\ref{sec:acet}章のモデルの要約で、各部分はそれぞれの章にもとづく；全体としては実験で検証されていない & 章内の導出 & 仮説 \\
12 & つまみは共通の制御なしに協調して働く：誤差、意外さ、書き込み、復帰（図~\ref{fig:timeline}） & 実験はない：模式図は定性的。個々の環は、\nt{NORA}と\nt{ACET}の分業の理論と、ヒトでの瞳孔と学習速度の関係 & \cite{YuDayan2005,Nassar2012} & 仮説 \\
13 & 一つの共通信号による学習は、結果に影響するニューロンが多いほど遅い（命題~\ref{prop:broadcastcost}） & 線形ネットワークの学習曲線の計算：出力の摂動による学習は勾配降下より出力の数の倍だけ遅い & \cite{Werfel2005} & 理論 \\
14 & 皮質が多層回路をどう調整するかは不明；候補はスパイクのバースト、ニューロンの枝の中での計算、予測による自己教師あり学習 & 誤差の担い手としてのバーストはモデル；皮質ニューロンの詳細モデルの応答を再現できたのは $5$--$8$ 層のネットワークだけ（モデル）；ヒト皮質ニューロンの枝で排他的論理和を生むカルシウムスパイクはスライスで測定済み；自己教師あり学習は証拠の総説 & \cite{Payeur2021,Beniaguev2021,Gidon2020,Keller2018} & 仮説 \\
\bottomrule
\end{longtable}}
